\documentclass[script,con]{debate}
\motion{AI的迅猛发展提升 / 降低了人类创作者存在的意义}
\stance{反方}{降低了}
\meta{赛制}{招新赛 3v3}
\meta{口径来源}{备赛文档 第十节 反方口径表}
\begin{document}
\debatetitle
\begin{thesisbox}[全场一句话立论]AI 让社会对多数创作者的需要绕开了人，留下来的人手里只剩修补：用 AI，作品里的你变少；不用 AI，需要你的人变少。\end{thesisbox}

\begin{kit}{1. 反一开篇立论}{套件 · 组装后 807–812 字 / 180 秒}
\assembly{开头（任选 1） → 立论主体}
\begin{pick}{开头}{1}
\begin{module}{如果质询拿到了正方承认「执行型订单在减少」}
谢谢主席。刚才质询里，正方已经承认，执行型的订单在减少。请评委记住这句承认，我方的第一点就从这里开始。
\end{module}
\begin{module}{如果质询里正方没有正面回答「被珍视却没人来请，算不算提升」}
谢谢主席。刚才有一个问题，正方没有正面回答：被珍视却没人来请，这算不算意义提升。这个问题，我方会一直留着。
\end{module}
\begin{fallback}{质询没有拿到明确成果}
谢谢主席，各位好。质询先放一边，我方先把今天到底在比什么说清楚，再请评委和我们一起算这笔账。
\end{fallback}
\end{pick}
\begin{fixed}{立论主体}
先说清三个词。AI 的迅猛发展，指的是 2022 年以来，机器生成文字、图像、音乐和声音的能力突然变强，而且可以不经过人，直接交出成品。人类创作者，是把创造文艺作品当成自己一部分的人，职业的、业余的都算；今天尤其要看那些不出名、靠作品被需要的大多数。存在的意义，词典里叫作用与重要性。放到今天是两个问题：这个世界还有多需要、多珍视由人来创作；创作还能给创作者自己留下多少。

所以判准是：和 AI 爆发之前比，由人来创作这件事，整体上变重了还是变轻了。需要和珍视两半都算，而且看多数创作者。我方要证明它变轻了；正方要证明它变重了，而且重在多数人身上，不只重在几个有名字的人身上。

我方第一点，不再被需要。大多数创作者不是大师，是画海报的插画师、写文案的写手、做字幕的译者、给动画配音的配音员。他们存在的意义里最实的一块，就是有人需要他们来画、来写、来说。现在，这份需要可以绕开他们了。美国圣路易斯华盛顿大学等机构的学者，拿一个大型自由职业平台做了对比：ChatGPT 发布之后，写作类自由职业者的月收入降了百分之五以上，做图像的降了将近一成。英国作家协会 2024 年的调查里，两成六的插画师、三成六的译者，已经因为生成式 AI 丢了活。说白了，被需要这一半，正在多数人身上缩水。

我方第二点，手艺被掏空。留下来的人，作品里属于他自己的那部分，正在变少。插画师 Amber Yu，曾经花整整一个星期，画一个穿着传统服装舞狮的女子，一张游戏海报收三千到七千元。后来游戏公司自己用 AI 几秒钟出图，只请她修修光线、修修画歪的手脚，价钱是原来的十分之一。你想想，画什么，甲方定了；怎么画，机器画了；她手里只剩修补。这不是一个人的遭遇。一项追踪五万多名用户、四百万件作品的研究发现，用了 AI 以后，作品的平均内容新颖度在下降；视觉上的新颖度，从最好的到平均的，全都在下降。\stress{作品越来越像，就是作品越来越不像某一个人。}

合起来，这是一个夹击：用 AI，作品里的你变少；不用 AI，需要你的人变少。所以我们全场都会请正方回答一个问题：你们说的被珍视，有多少落到了一个不出名的插画师身上？谢谢。
\end{fixed}
\end{kit}
\begin{contingency}
\ifthen{正方把意义定义成“只有人能给的东西”}{那就等于说 AI 替代什么都不算损失，这是把剩下的当成全部；词典里的意义是作用与重要性，被需要也算。}
\ifthen{正方说“订单只降了百分之二”}{那是订单数，收入降得更多；而且过去评价越好的人受冲击越大，交给二辩申论展开。}
\end{contingency}

\begin{stage}{2. 反一被质询防守预案}{正二质询，90 秒双边计时}
\textbf{原则}：先给是或否，再一句话守住：两半都算，但珍视落在少数，需要在多数身上缩水。不否认法律和标注是真的，否认常识会失去评委。
\begin{qa}
\qarow{贵方说的意义，包不包括被珍视？}{包括，两半都算。但珍视落在少数有名字的人身上，需要在多数人身上缩水。}
\qarow{法院只认人当作者，算不算看重人？}{算法律上的区分。它拦住了机器署名，没有给任何一个画师多带来一单生意。}
\qarow{珍视在变重，贵方还说整体变轻，对吗？}{对。变重的珍视落在少数人身上，多数人被需要的那一半在缩水，净额是变轻。}
\qarow{北京那个 AI 图的案子，判给人还是 AI？}{判给人。可法院认的独创性，已经缩到了提示词和参数。}
\qarow{法院看的，是不是人的选择和安排？}{是。选择在会写提示词的人手里，不一定在画了一星期海报的画师手里。}
\qarow{写作类月订单只降了百分之二吧？}{订单降百分之二，收入降百分之五以上，图像类收入降将近一成；而且过去评价越好的人受冲击越大。}
\qarow{到 2028 年那组数字是预测吧？}{是预测，我方标明了。已经发生的，看平台研究和作者调查，方向一致。}
\end{qa}
\end{stage}

\begin{stage}{3. 反二质询正一问题链}{90 秒，双边计时}
\textbf{总目标}：给反一立论的开头取材。拿到“执行型订单在减少”的承认，或者让“被珍视却没人来请”这一问留在场上。两链共五问。
\begin{chain}{链一：需要链}{让对方承认被需要的那一半在变轻}
\q{插画师因为 AI 丢了活，贵方承认吗？}{答“承认”：下一问；答“不承认”：追“英国两成六的插画师说丢了活，贵方否认？”}
\q{被需要，算不算存在意义的一部分？}{答“算”：下一问；答“不算”：点出对方把意义和被需要切开}
\q{那这一部分，是变重了还是变轻了？}{答“变轻”：收束；答“变重”：追“丢了活的插画师，是更被需要了吗？”}
\cut{好，我换个问法。}
\close{感谢，正方承认被需要的那一半在变轻。}
\end{chain}
\begin{chain}{链二：标签链}{把“被珍视却没人来请”留在场上}
\q{给作品标上“人做的”，能带来订单吗？}{答“能”：追“哪一单？”；答“不一定”：下一问}
\q{被珍视却没人来请，算意义提升吗？}{答“算”：收束，请评委记住；答非所问：复述一次，再不答就收束}
\close{感谢。这个问题，我方全场都会留着。}
\end{chain}
\end{stage}

\begin{stage}{4. 反三对辩要点}{双方各 90 秒}
\textbf{任务}：接住正三开出的第一个战场，把“需要和珍视”打成“珍视落在谁身上”。每次发言先一句回应，再抛一问。
\begin{points}{进攻点（4–5 个，每个 15–20 秒）}
\item 两半都算。可正方的珍视落在谁身上？一个不出名的插画师拿到了吗？
\item 标签保护的是不被冒名，还是有人来请？
\item 贴标签的那个实验，画全是 AI 画的：人们看重的是标签，还是人？
\item 八国九千人的盲听里，九成七分不出人和 AI。认不出来的珍视，落得到谁身上？
\item 用 AI，作品里的你变少；不用 AI，需要你的人变少。正方站哪一头？
\end{points}
\begin{points}{备用点（6–8 个）}
\item 对方说订单只降百分之二 → 收入降百分之五以上，图像类将近一成；两成六的插画师已经丢了活。
\item 对方说意义不等于订单 → 对不出名的创作者，被需要就是意义最实的一块。
\item 对方说法院只认人当作者 → 机器不能署名，不等于有人来请你。
\item 对方说标识 → 标识办法写明的目的是提醒用户辨别虚假信息。
\item 对方说更多人能创作 → 变多的是作品，不是被需要的作者。
\item 对方说获赞高了一半 → 那是主动用 AI 的网站用户，而且平均新颖度在降。
\item 对方说那是预测 → 是预测，我方标明；已经发生的看平台研究和作者调查。
\end{points}
\begin{points}{对方最可能问的三个问题 → 回应}
\item 意义包括被珍视吗？ → 包括，两半都算；但珍视落在少数，需要在多数身上缩水。
\item 八成的人要求标注，不是在乎人吗？ → 在乎的是别被骗。想知道外卖是不是预制菜，不等于会多付钱请厨师。
\item 月订单只降百分之二，能叫变轻？ → 那是订单数；收入降得更多，而且过去评价越好的人受冲击越大。
\end{points}
\end{stage}

\begin{kit}{5. 反二申论}{套件 · 组装后 381–393 字 / 90 秒}
\assembly{开头 → 拆正方立论（任选 1） → 推进论点一 → 收尾}
\begin{fixed}{开头}
谢谢主席，各位好。这是全场第一篇申论，我做两件事：先看正方立论的地基稳不稳，再把我方第一点往前推一步。
\end{fixed}
\begin{pick}{拆正方立论}{1}
\begin{module}{如果正方立论建立在「人味稀缺，所以更被珍视」上}
正方立论的地基，是人味稀缺，所以更被珍视。我方看三处。第一，被区分不等于被需要，标签保护的是不被冒名，不是有人来请。第二，最常被引的那个贴标签实验，画全是 AI 画的，人们看重的是一张标签；而八国九千人的盲听调查里，九成七的人分不出人和机器。第三，稀缺品的溢价，从来只流向能被认出名字的少数人，经济学家罗森早就讲过超级明星效应。所以正方要回答：珍视落在了谁身上？
\end{module}
\begin{module}{如果正方立论建立在「立意回到人」上}
正方立论的地基，是立意回到人。我方问一句：回到了哪个人？游戏公司自己用 AI 出图，立意在甲方手里，画师只被请来修光线、修手指。构思和执行被拆开，构思去了甲方，执行给了机器，画师手里只剩修补。讲立意最常被引的那项五万多名用户的研究，同一篇里写着：平均内容新颖度在下降，视觉新颖度从最好的到平均的全在下降。立意如果真回来了，作品为什么越来越像？
\end{module}
\begin{module}{如果正方立论建立在「法律和标注只认人」上}
正方立论的地基，是法律和标注只认人。我方看两处。第一，美国法院说的是 AI 不能当作者，没有给任何一个画师多带来一单生意；我国的标识办法，写明是为了提醒用户辨别虚假信息，不是替创作者加分。第二，被法律单独保护，往往说明处境变差了。说白了，没有人会给不愁的东西立法；濒危的物种才进保护名录。把一道防护栏说成一块奖牌，是看反了方向，也看错了对象。
\end{module}
\begin{fallback}{正方立论的主干不在以上几条时}
不管正方立论用什么论据，请评委看它量的是谁。判准要看多数创作者。法院的判决、实验室的评分，这类证据量的是规则和标签；而多数创作者，是平台上接单的写手、画海报的插画师、做字幕的译者、给动画配音的配音员。他们被需要的那一半有没有缩水，才是这道题真正的分量所在。正方的证据再漂亮，也要回答这一句。谁在讲多数，谁在讲少数，评委一听就清楚。
\end{fallback}
\end{pick}
\begin{fixed}{推进论点一}
我方第一点再往前走一步。不再被需要，连最好的人也挡不住。那项自由职业平台研究发现，过去评价越好、要价越高的顶尖自由职业者，受冲击反而更大。国际作者和作曲者协会联合会的预测更直白：到 2028 年，给影视做配音和字幕的译者，五成六的收入面临被 AI 拿走的风险。这是预测，但方向和已经发生的一致。你想想，连最好的那批人都挡不住，剩下的大多数，靠什么被需要？
\end{fixed}
\begin{fixed}{收尾}
所以我问一句：正方说的珍视，有多少真正落到了一个不出名的插画师身上？谢谢。
\end{fixed}
\end{kit}

\begin{stage}{6. 反二被质询防守预案}{正三质询，90 秒双边计时}
\textbf{原则}：反二刚推进了“连最好的人也挡不住”，正三最可能压数据的量级和“要求标注是在乎人”。承认数字，守住方向；承认在乎，分开“在乎”和“请你”。
\begin{qa}
\qarow{写作类月订单降百分之二，对吗？}{对。月收入降百分之五以上，图像类收入降将近一成，而且过去评价越好的人受冲击越大。}
\qarow{这个幅度，能证明整体变轻吗？}{这是 ChatGPT 刚发布之后的短期数据，方向已经清楚；两成六的插画师已经丢了活。}
\qarow{八成受访者要求标注，贵方否认吗？}{不否认。}
\qarow{要求标注，是不是在乎背后是不是人？}{是在乎别被骗。在乎，和请你画，是两回事。}
\qarow{配音字幕那组数字是预测吧？}{是预测，出自代表创作者的机构，我方标明了；只用它说趋势。}
\qarow{顶尖者受冲击更大，不正说明量的是执行？}{顶尖者的活恰恰最有个人风格。连他们都挡不住，替代的就不只是执行。}
\end{qa}
\end{stage}

\begin{stage}{7. 反三质询正二问题链}{90 秒，双边计时}
\textbf{总目标}：让正二承认标识的目的是防骗，认不出来的珍视落不到创作者身上。四问。
\begin{chain}{链一：防骗链}{把标识从“珍视”拉回“防骗”}
\q{标识办法写的目的，是不是防虚假信息？}{答“是”：下一问；答“不只是”：追“答记者问里写的是不是辨别虚假信息？”}
\q{防骗的标签，能算对画师的珍视吗？}{答“能”：追“哪一个画师因此多接了一单？”；答“不能”：收束}
\cut{好，我换个问法。}
\close{感谢，标识是防骗，不是给画师加分。}
\end{chain}
\begin{chain}{链二：认出链}{让对方承认认不出来的珍视落不到人身上}
\q{九成七的人盲听分不出人和 AI，对吗？}{答“对”：下一问}
\q{认不出来的珍视，落得到创作者身上吗？}{答“落不到”：收束；答“所以要标注”：追“标注能让哪个画师多接一单？”}
\close{感谢。认不出来，珍视就落不到人身上。}
\end{chain}
\end{stage}

\begin{stage}{8. 反一对辩要点}{双方各 90 秒}
\textbf{任务}：在中段收拢，把评委拉回判准：两半都算，看多数；珍视落在少数，需要在多数身上缩水。
\begin{points}{进攻点（4–5 个，每个 15–20 秒）}
\item 两半都算：给我一个不出名的创作者更被需要的证据。
\item 看多数：法院、实验室、明星个例，这些量的是多数还是少数？
\item 被区分不等于被需要：一张标签能带来一单生意吗？
\item 用 AI，作品里的你变少；不用 AI，需要你的人变少。正方站哪一头？
\end{points}
\begin{points}{备用点（6–8 个）}
\item 对方说集中度下降 → 那是主动用 AI 的网站用户内部，不是职业画师。
\item 对方说重新发歌的歌手 → 有唱片公司支持的明星，个例说明不了多数。
\item 对方说立意回到人 → 回到了甲方，不是画师。
\item 对方说七成六想分辨 → 想分辨是怕被骗。
\item 对方说获赞高了一半 → 平均新颖度在降，作品越来越像。
\item 对方说意义不等于订单 → 对不出名的人，被需要是意义最实的一块。
\end{points}
\begin{points}{对方最可能问的三个问题 → 回应}
\item 贵方承认珍视变重了吗？ → 承认区分更清楚了，法律和标注是真的；但珍视落在少数，需要在多数身上缩水，净额是变轻。
\item 甲方算不算创作者？ → 出于己意创造的人都算；可那个画了一星期海报的人，现在连己意都用不上了。
\item 业余创作者不靠订单，意义也降了吗？ → 业余者要的是被读到、被认出来；九成七的人盲听分不出人和 AI，他们被认出的机会也在变小。
\end{points}
\end{stage}

\begin{kit}{9. 反三申论}{套件 · 组装后 382–401 字 / 90 秒}
\assembly{开头 → 回应正三（任选 1） → 处理正方最强点（任选 1） → 推进论点二 → 收尾}
\begin{fixed}{开头}
谢谢主席，各位好。我先回应正三，再处理正方最有力的一点，最后推进我方第二点。
\end{fixed}
\begin{pick}{回应正三}{1}
\begin{module}{如果正三申论主打「订单少了不等于意义轻了」}
正三刚才说，订单少了不等于意义轻了。可对一个不出名的创作者来说，被需要就是他意义里最实的一块。把这一块拿走，剩下的珍视落在法律条文和别人的名字上，他摸不着，也拿不回家。
\end{module}
\begin{module}{如果正三申论主打「修图的是手，立意还在人」}
正三刚才说，修图的是手，立意还在人。那我问一句，海报画什么，是谁定的？是甲方。立意确实还在人手里，只是已经不在画师手里了；画师手里剩下的，是修光线、修手指。
\end{module}
\begin{module}{如果正三申论主打「更多人能创作」}
正三刚才说，AI 让更多人能创作。我方不反对表达。谁都有说话的权利，可人人都能生成以后，“创作者”这个身份就越来越轻。多出来的是会生成的人，不是被需要的作者。
\end{module}
\begin{fallback}{正三申论的重点不在以上几条时}
先回到判准。今天比的是由人来创作这件事整体变重还是变轻，需要和珍视都算，而且看多数创作者，不看最红的几个，也不看一个动人的例外。我方就按这个标准往下说。
\end{fallback}
\end{pick}
\begin{pick}{处理正方最强点}{1}
\begin{module}{如果正方全场最有力的是「法律和标注只认人」}
正方最有力的一点，是法律和标注只认人，把人和机器分开了。我方承认，这是真的。但请评委看它护住的是什么：不被冒名。一个插画师不会因为 AI 图打上了标识，就多接到一单。标识办法的答记者问写得很清楚，是为了提醒用户辨别虚假信息。被单独保护，恰恰说明处境在变差；一道防护栏，不是一块奖牌。把它当成奖牌，就把处境看反了。
\end{module}
\begin{module}{如果正方全场最有力的是「贴上人类标签评价更高」}
正方最有力的一点，是贴上人类标签，评价就更高。我方承认，这类实验做得很干净。可最常被引的那个实验里，画全是 AI 画的，人们看重的是一张标签；而标签要先被认出来。八国九千人的盲听测试里，九成七的人分不出人和 AI。你想想，一份要靠标签才能认出来的珍视，标签一撕就没了。认不出来的珍视，落不到创作者身上。
\end{module}
\begin{module}{如果正方全场最有力的是「AI 让普通人也能创作」}
正方最有力的一点，是 AI 让普通人也能创作。我方承认，门槛确实降了。可讲这件事最常被引的那项写作实验，同一篇里写着：用了 AI 点子的故事，彼此越来越像。每个人都往前走了一步，走到的却是同一个地方。一个音乐平台说，今年六月高峰时，每天约九万首纯 AI 曲目涌进来，超过新上传的一半。变多的是作品，变少的是属于某一个人的东西。
\end{module}
\begin{fallback}{正方最有力的一点不在以上几条时}
不管正方最有力的是哪一点，道理很简单。请评委问两件事。它落在多数创作者身上，还是少数人身上？它说的是被需要，还是只是被区分？判准要的，是多数创作者被需要和被珍视加在一起变重了。只证明人和机器被分开、被标注，证明不了这一点，因为分开之后，多数创作者被需要的那一半有没有回来？正方还欠一个证据，这笔账还没有算。
\end{fallback}
\end{pick}
\begin{fixed}{推进论点二}
我方第二点再加一层。手艺被掏空，不只发生在接私活的人身上。重庆一家游戏美术工作室，十五个角色设计师裁掉了五个；留下的人说，两个人就能干原来十个人的活。广东一位游戏画师说，AI 让她们更高效，也更累。这不是创作更自由了。是流水线更快了。说到底，作品里的你在变少，累却一点没少。
\end{fixed}
\begin{fixed}{收尾}
所以我问一句：作品越来越像，属于作者自己的部分，是多了还是少了？谢谢。
\end{fixed}
\end{kit}

\begin{stage}{10. 反三被质询防守预案}{正一质询，90 秒双边计时}
\textbf{原则}：正一会为结辩取材，最想拿到“立意在人手里”和“好处不只给头部”。可以承认事实，但每一句都接回夹击：立意在人手里，不在创作者手里；好处在用 AI 的人里分散，不用 AI 的人订单在少。
\begin{qa}
\qarow{决定一张海报画什么的，是人还是 AI？}{是人，是甲方；不是那个画了一星期的画师。}
\qarow{那立意是不是始终在人手里？}{在人手里，不在创作者手里。题目问的是人类创作者。}
\qarow{出于己意去创造的人，算不算创作者？}{算。可一个画了一星期海报的人，现在连己意都用不上，只剩修补，这就是变轻。}
\qarow{获赞集中度下降，贵方否认吗？}{不否认。那是主动用 AI 的网站用户内部。}
\qarow{集中度下降，是不是好处不只给了头部？}{在用 AI 的人里是。可不用 AI 的人订单在少，这就是我方的夹击。}
\qarow{视觉新颖度下降，说的只是平均吧？}{不是，原文写的是最好的和平均的都在降。}
\end{qa}
\end{stage}

\begin{stage}{11. 反一质询正三问题链}{90 秒，双边计时}
\textbf{总目标}：让正三承认立意不在画师手里、作品越来越像。第三问只在正三举了个例时问。
\begin{chain}{链一：甲方链}{让对方承认立意回到的是甲方}
\q{海报的立意，现在是甲方出还是画师出？}{答“甲方”：下一问；答“都是人”：追“那画师手里还剩什么？”}
\q{画师只负责修光线，算立意回到画师吗？}{答“不算”：收束；答“算”：请评委判断修光线算不算立意}
\cut{好，我换个问法。}
\close{感谢，立意回到的是甲方，不是画师。}
\end{chain}
\begin{chain}{链二：趋同链}{让对方承认属于作者的部分在变少}
\q{同一研究里，视觉新颖度全面下降，对吗？}{答“对”：下一问；答“那是平均”：追“原文是不是最好的也在降？”}
\q{作品越来越像，属于作者的部分是不是变少了？}{答“少”：收束；答“有立意的人更显眼”：下一问}
\q{一个明星的个例，代表多数创作者吗？}{只在正三举了个例时问。答“不代表”：收束}
\close{感谢。作品越来越像，属于作者的部分在变少。}
\end{chain}
\end{stage}

\begin{stage}{12. 反二对辩要点}{双方各 90 秒}
\textbf{任务}：守住判准，把全场收成夹击。不开新战场，每一句都落回“看多数”。
\begin{points}{进攻点（4–5 个，每个 15–20 秒）}
\item 用 AI，作品里的你变少；不用 AI，需要你的人变少。正方站哪一头？
\item 过去评价越好的人受冲击越大：质量能换来被需要吗？
\item 珍视落在谁身上？给我一个不出名的创作者更被需要的证据。
\item 修光线、修手指，算不算立意回到画师？
\end{points}
\begin{points}{备用点（6–8 个）}
\item 对方说意义不只是订单 → 对，两半都算；可珍视落在少数，需要在多数身上缩水。
\item 对方说法院把图判给人 → 认的独创性已经缩到提示词和参数。
\item 对方说集中度下降 → 那是用 AI 的网站用户内部。
\item 对方说八成要求标注 → 在乎的是别被骗。
\item 对方说那是预测 → 是预测；平台研究和作者调查已经发生。
\item 对方说更多人能创作 → 变多的是作品，不是被需要的作者。
\end{points}
\begin{points}{对方最可能问的三个问题 → 回应}
\item 标注和法律不是珍视吗？ → 是区分；区分之后，没有一个画师因此多接一单。
\item 丢活的是执行，贵方为什么说意义降了？ → 对多数创作者，执行就是他被需要的方式，也是作品里留下他自己的方式。
\item 贵方是不是在反对 AI？ → 不是。我方只回答题目：对多数创作者，意义变轻了。
\end{points}
\end{stage}

\begin{kit}{13. 反一结辩}{套件 · 组装后 377–414 字 / 90 秒}
\assembly{归纳分歧（任选 1） → 处理最强点（任选 1） → 封路（任选 1） → 论点收束 → 价值升华 → 结尾}
\begin{pick}{归纳分歧}{1}
\begin{module}{如果全场主战场落在「需要和珍视哪个算数」}
谢谢主席。今天的主战场，落在需要和珍视哪个算数。我方的回答是两个都算，但要看落在谁身上：珍视落在少数有名字的人身上，需要却在多数人身上缩水。
\end{module}
\begin{module}{如果全场主战场落在「修图工还是立意者」}
谢谢主席。今天的主战场，落在修图工还是立意者。我方的回答是：立意确实还在人手里，只是已经不在创作者手里了；创作者手里剩下的，是修补。
\end{module}
\begin{fallback}{主战场不清楚时}
谢谢主席。回到判准：和 AI 爆发之前比，由人来创作这件事，对多数创作者整体变重还是变轻。需要和珍视都要算，不能只看最红的几个人。
\end{fallback}
\end{pick}
\begin{pick}{处理最强点}{1}
\begin{module}{如果正方全场最有力的是「法律和标注只认人」}
正方最有力的，是法律和标注只认人。我方承认，这是真的。但法律说的是机器不能署名，没说有人会来请你。被单独保护，说明处境在变差；一道防护栏，不是一块奖牌。标注之后，那个修图的画师并没有多接到一单。
\end{module}
\begin{module}{如果正方全场最有力的是「AI 让普通人也能创作」}
正方最有力的，是普通人也能创作了。我方承认，门槛确实降了。但人人都能生成，作品越来越像。门槛降下来，进门的人多了，屋里却越来越像；变多的是作品，变少的是属于某一个人的东西。
\end{module}
\begin{fallback}{正方最有力的一点不在以上几条时}
正方最有力的那一点，也请评委用判准去量：它落在多数创作者身上吗？它说的是被需要，还是只是被区分？我方不否认区分是真的，只是区分之后，多数人被需要的那一半并没有回来。
\end{fallback}
\end{pick}
\begin{pick}{封路}{1}
\begin{module}{如果正方全场主打「人味稀缺所以更珍贵」}
正方全场都在说，稀缺让人更珍贵。待会儿的结辩如果还走这条路，请各位问一句：珍贵落在了谁身上？是有名字的少数，还是那个修光线、修手指的插画师？稀缺的东西，从来只给认得出名字的人。
\end{module}
\begin{module}{如果正方全场主打「意义不等于订单」}
正方全场都在说，意义不等于订单。待会儿的结辩如果还走这条路，请各位问一句：对一个不出名的创作者，拿掉被需要，他还剩什么？意义不等于订单，可订单是他被需要的样子。
\end{module}
\begin{module}{如果正方全场主打「一个重新能创作的人」}
正方全场最动人的，是一个借 AI 重新能创作的人。待会儿的结辩如果还落在这样的人身上，请各位问一句：他代表多数创作者，还是最动人的少数？一个动人的例外，证明不了多数人的处境。
\end{module}
\begin{fallback}{正方全场路线不在以上几条时}
不管正方待会儿怎么收，请各位带着判准去听：它说的是多数创作者，还是少数人；是被需要，还是只是被区分。这两个问题答不上，结论就落不了地。
\end{fallback}
\end{pick}
\begin{fixed}{论点收束}
说到底，我方两点合起来是一个夹击：用 AI，作品里的你变少；不用 AI，需要你的人变少。两头都在变轻，中间没有一条路是变重的。
\end{fixed}
\begin{fixed}{价值升华}
我想起那个插画师。她曾经花整整一个星期，画一个穿着传统服装舞狮的女子。后来，她打开甲方发来的图，把画歪的手改正，把光线调亮，交稿，拿原来十分之一的钱。那张图上没有她的名字，也没有她的一笔。我们今天争的，不是机器会不会画画，是一个人还能不能在自己的作品里找到自己。
\end{fixed}
\begin{fixed}{结尾}
那张图很好看。只是，\stress{它已经不需要她了}。谢谢。
\end{fixed}
\end{kit}
\begin{contingency}
\ifthen{时间不够}{先把处理最强点换成兜底，再压封路；价值升华和结尾不删。}
\ifthen{正方全场路线说不清}{封路用兜底，只重申判准，不猜正一会说什么。}
\end{contingency}

\end{document}
